\section{Introduction}

常识知识通常不会在文本中被显式说明，却是人类理解语言和判断事件可行性的基础。例如，把一只火鸡放进冰箱通常是合理的，而把一头大象放进冰箱则违反物体大小和容器容量之间的常识约束。对于大型语言模型而言，这类样本表面词汇相近、句式相似，但现实可行性完全不同，因此可以检验模型是否真正使用了稳定的世界知识，而不是依赖浅层词汇模式。

\comve 常识验证与解释任务由 Wang et al.~\citep{wang2019sensemaking,wang2020semeval} 提出并用于 SemEval-2020 Task 4。其中子任务 A 要求模型在两个相近句子中判断哪一个违背常识，是一个清晰、可复现的二分类任务。传统预训练编码器如 RoBERTa~\citep{liu2019roberta} 可以通过监督微调获得较强表现；近年来，大语言模型也能够通过提示学习、少样本学习和显式推理完成该任务。然而，仅报告一次性准确率并不能说明模型判断是否可靠。模型可能在原始测试集上表现较好，却对句子顺序、提示措辞、示例选择或采样随机性高度敏感。

近期研究也指出，常识推理和大语言模型可靠性需要从更细粒度角度评估。针对 \comve 的大语言模型评估发现，模型在验证任务上可能取得较高分数，但解释和稳定推理仍存在不足~\citep{comve_llm_2025}。HellaSwag-Pro 通过更困难和更多变的常识推理样本强调了题面扰动的重要性~\citep{hellaswagpro2025}。BrainBench 进一步显示，即使是强模型，在面向人类直觉的常识测试中仍会出现不稳定与退化~\citep{brainbench2026}。同时，test-time scaling 和自一致性研究表明，推理时多次采样、候选筛选和验证器重排可以在不重新训练模型的情况下改善部分推理任务~\citep{wang2023selfconsistency,testtime2025}。不确定性估计研究也表明，sample consistency 等黑盒置信度指标可用于分析 LLM 输出可靠性~\citep{uncertainty2025}。

本文将课程项目定位为常识验证中的大语言模型可靠性评估。我们保留 RoBERTa 微调和 LLM 提示式方法作为基础对照，同时设计两个并行扩展方向。第一个方向评估鲁棒性与一致性，关注输入扰动后模型答案是否保持等价。第二个方向研究验证增强与不确定性感知推理，关注模型能否通过显式常识约束、候选验证和采样一致性给出更可靠的判断。本文不加入中文翻译对比，因为当前数据集和官方任务均为英文，额外机器翻译会引入新的翻译质量变量，反而削弱对输入扰动和推理设置的聚焦。

本文的主要贡献体现在三个方面。首先，我们在 \comve 子任务 A 上构建 RoBERTa 与 Qwen 系列 LLM 的统一预测文件、指标文件和结果目录，并系统评估顺序交换、提示改写、少样本示例变化和采样随机性带来的鲁棒性问题。其次，我们探索 constraint-first、generate-verify-rerank 和 consistency uncertainty 等轻量推理时增强方法，观察它们是否能在不重新训练模型的情况下提升可靠性。最后，我们通过错误分析说明模型在物理常识、语言歧义、过度推理、prompt 敏感和高置信错误上的典型失败模式。
